% BHU PYQ -- Manually transcribed & error-corrected
% Paper: BCF-122 : Fundamentals of Business Finance
% Exam: B.Com. (Hons.) FMM Semester II Examination, 2022-23

\documentclass[11pt,a4paper]{article}
\usepackage[a4paper,top=2.5cm,bottom=2.5cm,left=2.8cm,right=2.8cm,headheight=14pt]{geometry}
\usepackage{amsmath,amssymb}
\usepackage[T1]{fontenc}
\usepackage[utf8]{inputenc}
\usepackage{lmodern,microtype,fancyhdr,enumitem,hyperref,lastpage,parskip}

\hypersetup{
    colorlinks=true,
    urlcolor=black,
    linkcolor=black,
    pdftitle={BCF-122: Fundamentals of Business Finance -- BHU B.Com. (Hons.) FMM Sem II 2022-23}
}

\pagestyle{fancy}
\fancyhf{}
\renewcommand{\headrulewidth}{0.4pt}
\renewcommand{\footrulewidth}{0.4pt}
\fancyhead[L]{\small   \textit{Banaras Hindu University}}
\fancyhead[R]{\small   \textit{BCF-122: Fundamentals of Business Finance}}
\fancyfoot[C]{\small Page  \thepage\ of \pageref{LastPage}}


\newlist{parts}{enumerate}{2}
\setlist[parts,1]{label=(\alph*),leftmargin=*,itemsep=5pt,topsep=3pt}
\setlist[parts,2]{label=(\roman*),leftmargin=*,itemsep=3pt,topsep=2pt}

\newcommand{\pts}[1]{\hfill   \textit{[#1]}}


\begin{document}


\begin{center}
  {\large\bfseries BANARAS HINDU UNIVERSITY}\\[2pt]
  {\normalsize B.Com. (Hons.) FMM --- Semester II Examination, 2022-23}\\[6pt]
  \rule{0.8\linewidth}{0.5pt}\\[6pt]
  {\large\bfseries Paper No. BCF-122: Fundamentals of Business Finance}\\[4pt]
  {\normalsize Time: 3 Hours\hfill Maximum Marks: 70}
\end{center}

\rule{\linewidth}{0.4pt}

\smallskip



\noindent   \textit{Instructions: Answer all questions. All questions carry equal marks (14 marks each).}

\bigskip




\noindent   \textbf{Question 1.}
``Finance is the life blood of industry.'' Elucidate this statement with suitable illustration.\pts{14}

\centerline{   \textbf{OR}}

``The goal of profit maximization does not provide an operationally useful criterion for measuring the success of business operations.'' Explain the above statement in context of the basic objective of the business finance.\pts{14}

\medskip\rule{\linewidth}{0.2pt}




\noindent   \textbf{Question 2.}
What are the different types of sources of finance? Explain long-term sources of finance.\pts{14}

\centerline{   \textbf{OR}}

What is Financial Planning? State its nature and types.\pts{14}

\medskip\rule{\linewidth}{0.2pt}




\noindent   \textbf{Question 3.}
``Inadequate capital is disastrous, whereas excess capital is a criminal waste.'' Critically evaluate.\pts{14}

\centerline{   \textbf{OR}}

``Investment, Financing and Dividend decisions are all inter-related.'' Comment.\pts{14}

\medskip\rule{\linewidth}{0.2pt}




\noindent   \textbf{Question 4.}
What is `Capital Budgeting'? Describe the steps involved in the project evaluation of business.\pts{14}

\centerline{   \textbf{OR}}

Define overall cost of capital. How will you determine the cost of capital from different sources?\pts{14}

\medskip\rule{\linewidth}{0.2pt}




\noindent   \textbf{Question 5.}
What are the essentials of an adequate dividend policy? Explain the various types of dividend.\pts{14}

\centerline{   \textbf{OR}}

Write short notes on any two of the following:\pts{14}

\begin{parts}
  \item Factors influencing capital structure of a company
  \item Stable dividend policy
  \item Weighted average cost of capital
\end{parts}

\vfill
\rule{\linewidth}{0.4pt}

\begin{center}\small
\href{https://bhu-exam.vercel.app/}{Click here to visit our website}\\[4pt]
\href{https://me-aryan.vercel.app/}{Click here to visit our contributor}\\[4pt]
\href{https://quashan-me.vercel.app/}{Click here to visit our contributor}
\end{center}

\end{document}
